\section{评价方法：答案正确之外还应测量什么}
\label{sec:evaluation}

\subsection{分层评价检索、生成与系统开销}

RAG 是组合系统，端到端准确率无法单独说明改进来自何处。检索侧需要判断必要证据是否被找到，以及上下文中混入了多少无关内容；生成侧需要判断答案是否正确、是否受到证据支持、是否妥善利用参数知识；系统侧则应记录检索次数、输入输出 Token、墙钟延迟和外部调用成本。对于自适应方法，还应有触发准确率、平均路径长度和无效循环比例等行为指标。

RAGAS 提出无参考答案场景下的自动评价，关注上下文相关性、答案忠实性和答案相关性，便于在开发过程中快速比较流程~\citep{es2024ragas}。ARES 使用合成训练数据训练轻量评价模型，并结合少量人工标注与预测驱动推断降低评价误差，其维度同样覆盖上下文相关性、答案忠实性和答案相关性~\citep{saadfalcon2024ares}。两者使模块化自动评价更可行，但评价器本身仍可能受到领域迁移、模型偏差和长答案分解错误影响，不能因此完全取消人工抽查。

RAGChecker 将参考答案和模型输出分解为细粒度声明，从整体、检索器和生成器三个层面提供诊断指标，例如上下文精度、声明召回、忠实性、幻觉和噪声敏感性~\citep{ru2024ragchecker}。这种声明级分析尤其适合自适应 RAG：如果增加检索后声明召回上升但噪声敏感性也上升，就能理解为何最终答案没有同步改善。相比只报告一个综合分数，它更接近对故障链的定位。

\subsection{基准不仅要有正常问题，还要有失败情境}

RGB 从噪声鲁棒性、拒答能力、信息整合和反事实鲁棒性四方面评价模型，为分析检索上下文的不同异常提供了中英双语基准~\citep{chen2024rgb}。其中，拒答和反事实情境尤其重要：若知识库不包含答案，系统应承认证据不足，而不是强迫检索器寻找一个看似相关的段落；若检索内容与模型记忆冲突，也不能只以“服从上下文”或“坚持参数知识”概括所有情况，而应检查任务设定与证据可信度。

CRAG Benchmark 则面向更真实和动态的事实问答，覆盖五个领域、八类问题，并通过模拟 API 提供网页与知识图谱搜索环境，纳入实体流行度和时间动态等变化~\citep{yang2024cragbench}。这类设置提醒研究者，静态维基百科上的单次问答不能代表真实 RAG 的全部难度。自适应策略在新闻、长尾实体或快速变化事实上的路由与停止行为，可能与标准数据集明显不同。

\begin{table}[t]
\centering
\small
\caption{自适应 RAG 的分层评价框架}
\label{tab:evaluation}
\begin{tabularx}{\textwidth}{p{2cm}p{4.1cm}XX}
\toprule
层级 & 核心问题 & 可用指标或测试 & 容易出现的误判 \\
\midrule
检索层 & 必要证据是否齐全，噪声是否可控 & 声明/证据召回、上下文精度、排序与去噪测试 & 相关文档多不等于证据覆盖完整 \\
生成层 & 答案是否正确且受证据支持 & 正确率、忠实性、引用支持度、幻觉率 & 流畅或忠实于错误文档不等于事实正确 \\
控制层 & 检索与停止是否发生在合适时机 & 不必要检索率、漏检率、平均轮数、无效循环率 & 只测最终答案会掩盖大量无效行动 \\
系统层 & 改进是否值得其资源开销 & 延迟、Token、检索/API 成本、吞吐 & 不同预算下的结果不能直接排名 \\
风险层 & 无答案、冲突与动态知识下是否可靠 & 拒答、校准、冲突证据和时效性测试 & 平均准确率会掩盖高风险失败 \\
\bottomrule
\end{tabularx}
\end{table}

\subsection{面向自适应方法的公平比较原则}

不同论文的结果通常来自不同检索器、知识库快照、生成模型和提示词，不能把表中最高准确率直接理解为控制策略普遍更优。本文认为，比较自适应 RAG 至少应遵守四项原则。

第一，设置共同基线，包括无检索、固定单次检索和固定多轮检索，并保持知识库与生成模型一致。第二，同时报告质量与成本曲线，而非只给单一操作点；可以通过调整触发阈值或最大轮数，观察在相同平均检索次数或延迟预算下谁的表现更好。第三，混入无需检索、无法回答和含噪声证据的问题，否则“总是检索”不会受到应有惩罚。第四，对控制器做消融，例如把学习到的触发器替换为随机或固定阈值策略，以判断收益究竟来自自适应决策还是更多调用。

自动评价还应保留可复核的证据链。对于每个答案，记录实际查询、返回文档、选择的行动及停止原因，再抽样检查评价器对“支持”和“冲突”的判断。尤其在评价模型与生成模型相同或同源时，二者共享偏差可能使分数偏高。因而，LLM-as-a-Judge 更适合作为规模化评价工具，而不是免于人工验证的最终裁判。
